
\documentclass[10pt,a4paper]{article}
\usepackage[top=12mm,bottom=12mm,left=14mm,right=14mm,headsep=2mm,footskip=7mm]{geometry}
\usepackage{fontspec}
\usepackage{xeCJK}
\setmainfont{Noto Sans}
\setsansfont{Noto Sans}
\setCJKmainfont{Noto Sans CJK SC}
\setCJKsansfont{Noto Sans CJK SC}
\usepackage{xcolor}
\usepackage{graphicx}
\usepackage{tikz}
\usetikzlibrary{calc,arrows.meta,positioning}
\usepackage{fancyhdr}
\usepackage{hyperref}
\usepackage{ragged2e}
\usepackage{microtype}
\definecolor{P2Paper}{HTML}{FCFAF7}
\definecolor{P2Ink}{HTML}{2D3238}
\definecolor{P2Muted}{HTML}{70757A}
\definecolor{P2BlueLight}{HTML}{D3EAF5}
\definecolor{P2Apricot}{HTML}{F5D9B8}
\definecolor{P2Rose}{HTML}{E8C7D3}
\definecolor{P2Violet}{HTML}{D7D3EA}
\definecolor{P2Blue}{HTML}{5B9FBE}
\definecolor{P2ApricotAccent}{HTML}{D49757}
\definecolor{P2RoseAccent}{HTML}{B97589}
\definecolor{P2VioletAccent}{HTML}{8177A8}
\definecolor{P2Rule}{HTML}{DDE3E7}
\definecolor{P2Panel}{HTML}{FFFDFC}
\pagecolor{P2Paper}
\color{P2Ink}
\setlength{\parindent}{0pt}
\setlength{\parskip}{3.5pt}
\linespread{1.04}
\hypersetup{hidelinks}
\pagestyle{fancy}
\fancyhf{}
\renewcommand{\headrulewidth}{0pt}
\fancyfoot[L]{\fontsize{6.6}{8}\selectfont\color{P2Muted}PROCESS REPORT · WORKFLOW CONSTRUCTION}
\fancyfoot[R]{\fontsize{6.6}{8}\selectfont\color{P2Muted}\thepage}
\newcommand{\eyebrow}[1]{{\fontsize{7.8}{9.4}\selectfont\bfseries\color{P2RoseAccent}#1}\par\vspace{1pt}}
\newcommand{\pagetitle}[1]{{\fontsize{18.5}{21.5}\selectfont\bfseries\color{P2Ink}#1}\par\vspace{3pt}}
\newcommand{\deck}[1]{{\fontsize{9.0}{13.0}\selectfont\color{P2Ink}\RaggedRight #1}\par\vspace{4pt}}
\newcommand{\tracehead}[2]{{\fontsize{7.1}{8.4}\selectfont\bfseries\color{#1}#2}\par}
\newcommand{\tracetitle}[1]{{\fontsize{10.2}{12.2}\selectfont\bfseries\color{P2Ink}#1}\par\vspace{1.5pt}}
\newcommand{\tracebody}[1]{{\fontsize{8.0}{11.2}\selectfont\color{P2Ink}\RaggedRight #1}\par}
\newcommand{\provenance}[1]{\vspace{2pt}{\color{P2Rule}\hrule height .35pt}\vspace{2.5pt}{\fontsize{6.9}{9.2}\selectfont\color{P2Muted}\RaggedRight\textbf{Original Evidence.} #1}\par}
\tikzset{
 userA/.style={rounded corners=1mm,draw=P2ApricotAccent,line width=.55pt,fill=P2Apricot,fill opacity=.34},
 userR/.style={rounded corners=1mm,draw=P2RoseAccent,line width=.55pt,fill=P2Rose,fill opacity=.34},
 gptB/.style={rounded corners=1mm,draw=P2Blue,line width=.55pt,fill=P2BlueLight,fill opacity=.27},
 gptV/.style={rounded corners=1mm,draw=P2VioletAccent,line width=.55pt,fill=P2Violet,fill opacity=.27},
 arrA/.style={-{Stealth[length=2.0mm,width=1.45mm]},draw=P2RoseAccent,line width=.72pt},
 arrB/.style={-{Stealth[length=2.0mm,width=1.45mm]},draw=P2Blue,line width=.72pt},
 arrV/.style={-{Stealth[length=2.0mm,width=1.45mm]},draw=P2VioletAccent,line width=.72pt},
 tagA/.style={circle,minimum size=4.8mm,inner sep=0pt,draw=P2ApricotAccent,line width=.6pt,fill=P2Paper,font=\fontsize{6.7}{7}\selectfont\bfseries},
 tagR/.style={circle,minimum size=4.8mm,inner sep=0pt,draw=P2RoseAccent,line width=.6pt,fill=P2Paper,font=\fontsize{6.7}{7}\selectfont\bfseries},
 tagB/.style={circle,minimum size=4.8mm,inner sep=0pt,draw=P2Blue,line width=.6pt,fill=P2Paper,font=\fontsize{6.7}{7}\selectfont\bfseries}
}
\begin{document}

\eyebrow{WORKFLOW CONSTRUCTION · 17}
\pagetitle{从“Codex说完成”到可复现工程：GitHub成为验收依据}
\deck{早期我已经把Codex定位成工程执行者。到了智慧城市项目，我又加了一层要求：重要阶段做完以后先push到GitHub，再让GPT直接读真实远程代码、结果和报告审核，不能只听Codex说“完成了”。}
\begin{tikzpicture}[x=1mm,y=1mm]
  \node[anchor=north west,inner sep=0] (u) at (0,0) {\includegraphics[width=45mm]{assets/user_git.png}};
  \node[anchor=north west,inner sep=0] (g) at (0,-38) {\includegraphics[width=72mm]{assets/gpt_git_workflow.png}};
  \begin{scope}[shift={(u.south west)},x={($(u.south east)-(u.south west)$)},y={($(u.north west)-(u.south west)$)}]
    \path[userA] (0.03,0.55) rectangle (0.98,0.80); \coordinate (ugit) at (0.98,0.675);
  \end{scope}
  \begin{scope}[shift={(g.south west)},x={($(g.south east)-(g.south west)$)},y={($(g.north west)-(g.south west)$)}]
    \path[gptB] (0.18,0.37) rectangle (0.98,0.68); \coordinate (gflow) at (0.98,0.525);
    \path[gptV] (0.18,0.33) rectangle (0.80,0.37); \coordinate (gaudit) at (0.80,0.35);
  \end{scope}
  \draw[arrA] (ugit) .. controls (80,-4) and (80,-61) .. (gflow);
  \draw[arrB] (ugit) .. controls (85,-7) and (85,-83) .. (gaudit);
  \node[tagA] at ($(ugit)+(3.8mm,0)$) {1};
  \node[anchor=north west,text width=87mm,align=left] at (92,-2) {
    \tracehead{P2ApricotAccent}{MICRO TRACE 19}\tracetitle{“给我Codex Prompt，并由你检查返回结果”}
    \tracebody{我最初只要求GPT帮我分析任务、写Codex Prompt并检查返回结果。后来为了让检查有真实依据，我把“检查结果”进一步改成了直接检查GitHub远程仓库。}
    \vspace{7mm}
    \tracehead{P2Blue}{ENGINEERING LOOP}\tracetitle{实现/运行 $\rightarrow$ push $\rightarrow$ GPT remote audit $\rightarrow$ 修复}
    \tracebody{这样GitHub就不只是最后存代码的地方，而是整个过程中用来验收的真实依据。Codex自报PASS之后，GPT仍然要看真实代码、结果、图表和报告。}
  };
\end{tikzpicture}
\provenance{这条要求来自我早期的真实协作需求；GPT截图来自智慧城市项目中后来整理出的GitHub协作流程。}
\newpage
\eyebrow{WORKFLOW CONSTRUCTION · 18}
\pagetitle{完整工程镜像与老师提交包必须分开}
\deck{这里我还特意把两件事分开：GitHub尽量保留完整工程，方便复现和审核；最后交给老师的Submission Package再按要求精简。不能因为老师最后不交某个文件，就提前把它从工程里删掉。}
\begin{tikzpicture}[x=1mm,y=1mm]
  \node[anchor=north west,inner sep=0] (g) at (0,0) {\includegraphics[width=75mm]{assets/gpt_fullmirror_submission.png}};
  \begin{scope}[shift={(g.south west)},x={($(g.south east)-(g.south west)$)},y={($(g.north west)-(g.south west)$)}]
    \path[gptB] (0.18,0.46) rectangle (0.98,0.67); \coordinate (gfull) at (0.98,0.565);
    \path[gptV] (0.18,0.16) rectangle (0.98,0.43); \coordinate (gsub) at (0.98,0.295);
  \end{scope}
  \node[anchor=north west,text width=88mm,align=left] at (89,-2) {
    \tracehead{P2Blue}{FULL MIRROR}\tracetitle{GitHub保存完整项目工程}
    \tracebody{代码、Notebook、实验结果、报告、图表、evidence和必要历史成果都可以保留在GitHub里。这样后面无论是我、GPT还是老师，都能追到结果是怎么产生的。}
    \vspace{8mm}
    \tracehead{P2VioletAccent}{SUBMISSION PACKAGE}\tracetitle{老师提交包只做最后一层整理}
    \tracebody{最终ZIP不需要把整个仓库原样复制一遍，只放老师明确要求和必要的复现文件。内部Prompt、日志、requirement matrix等可以留在工程里，不必都交给老师。}
    \vspace{8mm}
    \tracehead{P2RoseAccent}{WHY THIS MATTERS}\tracetitle{“完整”与“提交简洁”不再互相冲突}
    \tracebody{这样“工程要完整”和“提交包要干净”就不冲突了：GitHub负责真实性和复现，Submission Package负责最后提交。}
  };
\end{tikzpicture}
\provenance{本页使用真实讨论中“GitHub完整工程 vs Teacher Submission Package”的连续规则，不展示后期内部Upload Bundle整理过程。}
\end{document}
